\documentclass[a4paper,12pt]{article}

% Язык и кодировка
\usepackage[T2A]{fontenc}
\usepackage[utf8]{inputenc}
\usepackage[russian]{babel}

% Поля страницы
\usepackage[a4paper,top=2cm,bottom=2cm,left=3cm,right=2cm]{geometry}

% Полезные пакеты
\usepackage{amsmath}
\usepackage{graphicx}
\usepackage{array}
\usepackage{xcolor}
\usepackage{listings}
\usepackage[colorlinks=true, allcolors=blue]{hyperref}

% Оформление кода
\lstset{
  language=Python,
  basicstyle=\ttfamily\small,
  keywordstyle=\color{blue},
  commentstyle=\color{gray},
  frame=single,
  breaklines=true,
  showstringspaces=false
}

\title{Лабораторная работа №2\\Документирование}
\author{Студент Cейфуллаев Рамил Рахимович, ИСУ 561617}
\date{4 октября 2026}

\begin{document}
\maketitle

\begin{abstract}
В работе документируется файл \texttt{mysterious\_file.py}, оставленный предыдущим разработчиком. Функции получили осмысленные имена и docstring с примерами вызова, три функции с неочевидным поведением исправлены. Создан каталог \texttt{docs} с документацией в формате Markdown и история изменений ветки с хешами коммитов.
\end{abstract}

\section{Цель работы}

Освоить документирование кода: писать docstring, давать функциям и переменным осмысленные имена, оформлять отдельную документацию в Markdown и фиксировать историю изменений в Git.

\section{Ход работы}

\begin{enumerate}
\item Репозиторий \texttt{itmo\_tools\_2026}, переход на ветку \texttt{main\_features\_561617}.
\item Создание ветки командой \texttt{git checkout -b doc\_561617}.
\item Разбор функций в \texttt{mysterious\_file.py}: определение логики, типов параметров и результата.
\item Переименование функций и переменных, добавление docstring с описанием, \texttt{Args}, \texttt{Returns} и примером вызова.
\item Проверка: запуск файла и проверка примеров командой \texttt{python3 -m doctest mysterious\_file.py}.
\item Создание каталога \texttt{docs} с файлом \texttt{documentation.md}.
\item Слияние ветки \texttt{doc\_561617} в \texttt{main\_features\_561617} и отправка обеих веток на GitHub.
\end{enumerate}

\section{Описание функций}

\begin{table}[h]
\centering
\begin{tabular}{|l|l|p{5.5cm}|}
\hline
Было & Стало & Что делает \\ \hline
\texttt{buy\_products} & \texttt{calculate\_change} & считает сдачу после покупки \\ \hline
\texttt{check} & \texttt{is\_weekend} & проверяет, выходной ли день \\ \hline
\texttt{h} & \texttt{is\_even} & проверяет чётность числа \\ \hline
\texttt{z} & \texttt{multiples\_of\_three\_in\_range} & числа, кратные 3, в диапазоне \\ \hline
\texttt{sentiment\_score} & \texttt{sentiment\_score} & оценка настроения текста \\ \hline
\texttt{f} & \texttt{to\_binary\_string} & число в двоичную строку \\ \hline
\end{tabular}
\caption{Переименование функций}
\label{tab:functions}
\end{table}


\subsection{calculate\_change}

Бывшая \texttt{buy\_products}. Считает сдачу после покупки нескольких товаров. Если денег хватает, возвращает сдачу, иначе строку с сообщением о нехватке денег, поэтому тип результата зависит от условия. Логика не менялась, изменены только имя и документация.

\begin{lstlisting}
# Before
def buy_products(money, price_per_item, items_count):
    total = price_per_item * items_count
    if total > money:
        return "Недостаточно денег"
    return money - total

# After
def calculate_change(money, price_per_item, items_count):
    total = price_per_item * items_count
    if total > money:
        return "Недостаточно денег"
    return money - total
\end{lstlisting}

Функции \texttt{is\_weekend} и \texttt{is\_even} (бывшие \texttt{check} и \texttt{h}) также работали корректно, их логика не менялась.


\section{Исправленные функции}

\subsection{multiples\_of\_three\_in\_range}

В исходной функции при шаге больше 1 терялась часть кратных трём чисел: для диапазона $(1, 10)$ с шагом 2 возвращалось \texttt{[3, 9]} без числа 6. Параметр шага удалён.

\begin{lstlisting}
# Before
def z(a, b, c=1):
    r = []
    for i in range(a, b, c):
        if i % 3 == 0:
            r.append(i)
    return r

# After
def multiples_of_three_in_range(start, stop):
    result = []
    for i in range(start, stop):
        if i % 3 == 0:
            result.append(i)
    return result
\end{lstlisting}

\subsection{sentiment\_score}

Знаки препинания не отбрасывались, поэтому слово \texttt{love!} не совпадало со словом \texttt{love} из списка и баллов не давало. Теперь при поиске в списках знаки по краям слова убираются.

\begin{lstlisting}
# Before
if word in p_w:
    s += 1
elif word in n_w:
    s -= 1

# After
clean_word = word.strip('!?.,;:')
if clean_word in positive_words:
    score += 1
elif clean_word in negative_words:
    score -= 1
\end{lstlisting}

\subsection{to\_binary\_string}

Исходная функция собирала биты от младшего к старшему и разворачивала результат только при счётчике больше 10, поэтому \texttt{f(10)} возвращала \texttt{0101} вместо \texttt{1010}. Счётчик удалён, разворот выполняется всегда.

\begin{lstlisting}
# Before
def f(x, y=0):
    z = []
    while x > 0:
        z.append(x % 2)
        x = x // 2
        y += 1
    if y > 10:
        z.reverse()
    return ''.join(map(str, z)) if z else "0"

# After
def to_binary_string(number):
    bits = []
    while number > 0:
        bits.append(number % 2)
        number = number // 2
    bits.reverse()
    return ''.join(map(str, bits)) if bits else "0"
\end{lstlisting}

\begin{table}[h]
\centering
\begin{tabular}{|l|l|l|}
\hline
Вызов & До исправления & После исправления \\ \hline
\texttt{f(10)} & \texttt{0101} & \texttt{1010} \\ \hline
\texttt{sentiment\_score("love!", ["love"], ["hate"])} & 0 & 1.5 \\ \hline
\end{tabular}
\caption{Результаты до и после исправлений}
\label{tab:results}
\end{table}

\section{История изменений ветки}

Ветка \texttt{doc\_561617} создана от \texttt{main\_features\_561617} (коммит \texttt{e8ff79c}).

\begin{table}[h]
\centering
\begin{tabular}{|l|p{9cm}|}
\hline
Хеш & Сообщение \\ \hline
\texttt{abffddf} & Декларирование и исправление функций в mysterious\_file.py \\ \hline
\texttt{f4251e0} & Добавлена markdown-документация в каталог docs \\ \hline
\end{tabular}
\caption{Коммиты ветки}
\label{tab:commits}
\end{table}

Слияние в \texttt{main\_features\_561617} прошло как fast-forward, без конфликтов.

\section{Вывод}

В ходе работы все функции из \texttt{mysterious\_file.py} получили осмысленные имена и docstring с примерами вызова. Выявлены и исправлены три функции с неочевидным поведением. Документация оформлена в каталоге \texttt{docs} в формате Markdown и в формате \LaTeX{}.

\end{document}